\documentclass[a4paper,11pt]{article}
\usepackage[utf8]{inputenc}
\usepackage[T1]{fontenc}
\usepackage[ngerman]{babel}
\usepackage{amsmath,amssymb}
\usepackage[table]{xcolor}
\usepackage{tikz}
\usepackage{geometry}
\usepackage{array}
\geometry{margin=2cm}

% ---- Dark-Mode-Farben (Pflicht) ----
\definecolor{bgdark}{HTML}{1E1E1E}
\definecolor{fgtext}{HTML}{E6E6E6}
\definecolor{accent}{HTML}{89B4FA}
\definecolor{accent2}{HTML}{F38BA8}
\definecolor{gridcol}{HTML}{4A4A4A}
\definecolor{panel}{HTML}{2A2A2A}

\pagecolor{bgdark}
\color{fgtext}
\arrayrulecolor{fgtext}

\setlength{\parindent}{0pt}

% ---- Karo-Raster ----
\newcommand{\karo}[1]{\par\nobreak\vspace{0.4em}\noindent
  \begin{tikzpicture}
    \draw[step=5mm, gridcol, very thin] (0,0) grid (\linewidth,#1);
  \end{tikzpicture}\par}

% ---- Teilaufgabe zusammenhalten ----
\newcommand{\sub}{\filbreak\medskip}

% ---- Aufgabenkopf ----
\newcommand{\aufgabe}[3]{\clearpage
  {\large\bfseries\color{accent} Aufgabe #1\quad #2}\hfill{\bfseries\color{fgtext} #3 Punkte}\par
  \vspace{0.3em}\noindent\textcolor{gridcol}{\rule{\linewidth}{0.8pt}}\par\vspace{0.8em}}

\begin{document}

% ================= KOPF =================
\begin{center}
  {\LARGE\bfseries\color{accent} Probeklausur II -- Analysis II}\\[0.3em]
  {\normalsize Prof. Dr. Daniel Blochinger \quad$\cdot$\quad DHBW Ravensburg}
\end{center}
\vspace{0.6em}
\noindent\textcolor{gridcol}{\rule{\linewidth}{1pt}}\par\vspace{0.8em}

\renewcommand{\arraystretch}{1.8}
\noindent
\begin{tabular}{@{}p{2.2cm}p{5.3cm}@{\quad}p{3.6cm}p{4.5cm}@{}}
\textbf{Name:} & \textcolor{gridcol}{\hrulefill} & \textbf{Matrikel-Nr.:} & \textcolor{gridcol}{\hrulefill}\\
\textbf{Datum:} & \textcolor{gridcol}{\hrulefill} & \textbf{Bearbeitungszeit:} & 90 Minuten\\
\end{tabular}
\renewcommand{\arraystretch}{1.0}

\vspace{0.6em}
\noindent\textbf{Gesamtpunktzahl:} 50 Punkte\\[0.4em]
\noindent\textbf{Zugelassene Hilfsmittel:} Stifte, Lineale \& Geodreiecke sowie gedruckte und
handgeschriebene Unterlagen (Skripte, Mathebücher, selbst erstellte Formelsammlungen usw.).\\[0.2em]
\noindent\textbf{Nicht zulässig:} Taschenrechner und sämtliche computerartigen und/oder
internetfähigen Geräte (Laptops, Tablets, E-Book-Reader, Smartphones, Smartwatches usw.).

\vspace{1.0em}

\renewcommand{\arraystretch}{1.6}
\noindent
\begin{tabular}{|>{\bfseries}l|c|c|c|c|c||c|}
\hline
Aufgabe & 1 & 2 & 3 & 4 & 5 & $\Sigma$\\\hline
Thema & Diff.rechnung & Vektoranalysis & DGL & Integralr. & Trigonometrie & \\\hline
erreicht & von 10 & von 15 & von 12 & von 8 & von 5 & von 50\\\hline
\end{tabular}
\renewcommand{\arraystretch}{1.0}

\vspace{1.0em}
\noindent\textit{Hinweis: Begründen Sie Ihre Rechenwege nachvollziehbar. Zwischenschritte werden
bewertet (Folgefehler werden berücksichtigt).}

% ================= AUFGABE 1 =================
\aufgabe{1}{Wiederholung Differenzialrechnung}{10}

\textbf{a)} Bestimmen Sie die erste Ableitung der folgenden Funktionen.\hfill\textbf{[4,5 Punkte]}
\vspace{0.3em}

\sub(i)\quad $f(x) = x^2\, e^{3x}$
\karo{3cm}

\sub(ii)\quad $g(x) = \cos(x)\,\sin(x)$
\karo{3cm}

\sub(iii)\quad $h(x) = \dfrac{2x+1}{x^2+1}$
\karo{3.2cm}

\sub\textbf{b)} Gegeben sei die Funktion
\[
  f(x) = \begin{cases} x^2 + 1 & \text{für } x \le 2,\\[2pt] a\,x + b & \text{für } x > 2.\end{cases}
\]
Bestimmen Sie $a, b \in \mathbb{R}$ so, dass $f$ an der Stelle $x=2$ stetig \emph{und}
differenzierbar ist.\hfill\textbf{[5,5 Punkte]}
\karo{6cm}

% ================= AUFGABE 2 =================
\aufgabe{2}{Vektoranalysis}{15}

\sub\textbf{a)} Gegeben sei das Skalarfeld $f(x,y) = x^2 + y^2$ sowie der Punkt $P=(3,4)$.\hfill\textbf{[6 Punkte]}
\begin{itemize}\setlength\itemsep{0pt}
  \item[(i)] Berechnen Sie den Gradienten $\operatorname{grad} f$ und werten Sie ihn in $P$ aus.
  \item[(ii)] Geben Sie den \emph{Betrag der höchsten Steigung} in $P$ und die \emph{Richtung des
        steilsten Anstiegs} an.
  \item[(iii)] Berechnen Sie die Richtungsableitung von $f$ in $P$ in Richtung $\vec v = (4,3)$
        (auf Länge $1$ normieren).
\end{itemize}
\karo{6.5cm}

\sub\textbf{b)} Bestimmen Sie das totale Differenzial $df$ von $f(x,y)=x\,y^2$ und schätzen Sie damit
$f(3{,}1\,;\,1{,}9)$ näherungsweise ab (ausgehend von $(x,y)=(3,2)$).\hfill\textbf{[3 Punkte]}
\karo{4.5cm}

\sub\textbf{c)} An welchen Stellen erfüllt $f(x,y) = x^2 + y^2 - xy - 3x$ die \emph{notwendige}
Bedingung für eine Extremstelle?\hfill\textbf{[2 Punkte]}
\karo{4.5cm}

\sub\textbf{d)} Gegeben seien die Vektorfelder
\[
  \vec f(x,y) = \begin{pmatrix} x^2 - y^2 \\[2pt] 2xy \end{pmatrix},
  \qquad
  \vec g(x,y,z) = \begin{pmatrix} z \\[2pt] x \\[2pt] y \end{pmatrix}.
\]
Berechnen Sie $\operatorname{div}\vec f$ und $\operatorname{rot}\vec f$ sowie $\operatorname{rot}\vec g$.
Geben Sie für $\vec f$ zusätzlich an, unter welchen Bedingungen Quellen- bzw. Wirbelfreiheit
besteht.\hfill\textbf{[4 Punkte]}
\karo{6.5cm}

% ================= AUFGABE 3 =================
\aufgabe{3}{Gewöhnliche Differenzialgleichungen}{12}

\sub\textbf{a)} Klassifizieren Sie die folgenden Differenzialgleichungen (linear/nichtlinear,
homogen/inhomogen bzw. getrennte Veränderliche). Lösen müssen Sie hier noch
\emph{nicht}.\hfill\textbf{[2 Punkte]}\\[0.3em]
(i)\ $y' = x\,y^2$ \qquad (ii)\ $y' = \tfrac{1}{x}\,y + x$ \qquad (iii)\ $y' = 2y$
\karo{3.5cm}

\sub Lösen Sie nun die folgenden Anfangswertprobleme.
\vspace{0.3em}

\sub\textbf{b)}\quad $y' = x^2\,y^2, \qquad y(0) = 1$\hfill\textbf{[3 Punkte]}
\karo{5cm}

\sub\textbf{c)}\quad $y' = y + e^{2x}, \qquad y(0) = 0$\hfill\textbf{[4 Punkte]}
\karo{6cm}

\sub\textbf{d)}\quad $y' = \dfrac{3}{x}\,y, \qquad y(1) = 2$\hfill\textbf{[3 Punkte]}
\karo{4.5cm}

% ================= AUFGABE 4 =================
\aufgabe{4}{Integralrechnung}{8}

Berechnen Sie die folgenden Integrale.
\vspace{0.3em}

\sub\textbf{a)}\quad $\displaystyle \int_1^2 \left(\frac{3}{x^2} + \frac{1}{x}\right) dx$\hfill\textbf{[2 Punkte]}
\karo{4cm}

\sub\textbf{b)}\quad $\displaystyle \int x\, e^{x}\, dx$ \quad(partielle Integration)\hfill\textbf{[2,5 Punkte]}
\karo{4cm}

\sub\textbf{c)}\quad $\displaystyle \int 2x\,\cos(x^2)\, dx$ \quad(Substitution)\hfill\textbf{[1,5 Punkte]}
\karo{3.5cm}

\sub\textbf{d)}\quad $\displaystyle \int_0^1 \int_0^1 (x + y)\; dx\, dy$\hfill\textbf{[2 Punkte]}
\karo{4cm}

% ================= AUFGABE 5 =================
\aufgabe{5}{Trigonometrie}{5}

\sub\textbf{a)} Rechnen Sie $210^\circ$ ins Bogenmaß um und bestimmen Sie $\sin$ und $\cos$ dieses
Winkels.\hfill\textbf{[2 Punkte]}
\karo{3.5cm}

\sub\textbf{b)} Eine Schwingung wird durch $s(t) = 5\sin(2\pi t)$ beschrieben. Bestimmen Sie
Amplitude, Kreisfrequenz $\omega$, Frequenz $f$ und Periodendauer $T$.\hfill\textbf{[1,5 Punkte]}
\karo{3cm}

\sub\textbf{c)} Differenzieren Sie:\hfill\textbf{[1,5 Punkte]}\\
(i)\ $u(x) = \sin(x)\cos(x)$\qquad (ii)\ $v(x) = \cos(3x)$
\karo{3.5cm}

\end{document}
